\subsection{有价值的表示暴露有价值的复用单位}
Cartesian码本允许把每份primary表按坐标精确拆成两张16项表。
新版扫描中，该选择在22/24个单元领先；四组与八组都只存每子空间16项，
却执行不同次数的查找（\cref{fig:7-lut}）。因此，表示让紧凑评分器成为可能，
而表大小不能单独决定最佳实现。真正复用的单位是查询相关坐标组：其距离
构造一次，再被多个候选使用。

共享内存可容纳性解释了重要的资源转变：所测配置中，完整表加扫描scratch
超过可用空间，两组表可以放下。继续缩表时，额外查找与累加开始竞争。
packed word同样在逻辑字节量不变时改变指令和寻址工作。因此应联合考虑
状态、指令与复用；没有计数器时，无法进一步区分cache、bank、寄存器和issue。
两个例外LUT单元及下文Mailman案例反驳普遍的“最小表最好”规则。
两组默认来自实测取舍，不是跨设备最优结论。

这种机会比任意PQ更窄：这里的乘积结构存在于每个primary子码本\emph{内部}。
学习式PQ子空间的质心通常相互耦合，拆标签位不等于拆距离。
反过来，为降低重建误差而改码本，也可能失去该执行优势。
表示应结合评分器和残差层次可达的recall--成本前沿评价，
这把第5章的数学兼容性检查连接到第7章的端到端问题。

\subsection{有效并行受重复状态与选择成本限制}
相对lane精排参考，集成调度在batch 8的吞吐几何平均为$2.727\times$，
batch 1,000只有$1.015\times$（\cref{fig:7-batch}）。小batch缺乏查询级并行时，
把单查询分给多block可以补足工作；额外组也引入局部缓冲、重复query／table
staging及合并。原grid已有足够block时，这些成本可吞掉收益。
因此最大occupancy本身不能作为优化目标。

负载偏斜提供另一种分组理由。E1大batch只有Vogue走分组，候选数CV为
0.338--0.622；0.97处相对RaBitQ的比值，从单block的0.933变为调度路径的
1.059（\cref{tab:7-e1-paths}）。配对运行把执行机制连接到竞争区间，
但测的是调度、保留、共享界与合并的整体，单独识别共享界需要自己的消融。
浅probe的部分Vogue设置会变慢，BGE batch 256的调度regret达到1.217。
CV描述不均衡，却没有纳入绝对工作量、staging和merge成本。

工程上应保留query-centric路径，并同时依据可用并行度与单查询预计工作选择分组。
候选数、survivor容量和局部缓冲成本是更丰富决策的候选输入。
拟合选择器还需留出工作负载与校准成本核算，目前观察未验证这类选择器。
现有启发式的价值也在于其收益与失败单元可以分别复查。

\subsection{执行变快会改变避免精排的价值}
降低预算省去残差读取、计算与最终选择，路由和广泛primary评分仍保留。
可省比例取决于这些操作原本占多少延迟。更快的lane精排可能降低减少候选的回报，
更快primary扫描则可能提高精排的相对占比。二者都不会改变丢弃特定候选的质量后果。

这是成本模型推论，不是对每项历史／新版预算差异的因果解释。
ADC旧校准时间和收益属于不同执行记录，同时改变多个条件。
因果测试应固定预算决策与候选集合，只替换精排内核。
选中臂QPS还独立于校准计时，真正服务节省需要足够复用，四单元直接轨迹现已测量该边界
（\cref{sec:6-5}）。

新版FULL在全部12个单元满足所测$10^{-3}$相对损失目标，提供实测质量优先工作点。
小样本降低校准成本，却可能接受有害预算。独立结果位模型的分辨率计算支持
约500条的样本尺度，较强独立性和排序条件仍明确保留，补证十二个FULL枚举均已一致。
工程决策需要依据质量要求和预计复用来投入校准工作，单独优化选择器或只报告
有利的平均选中臂吞吐都不足以完成该决策。
